%% LyX 2.4.4 created this file.  For more info, see https://www.lyx.org/.
%% Do not edit unless you really know what you are doing.
\documentclass[oneside,english]{amsart}
\usepackage[T1]{fontenc}
\usepackage[latin9]{inputenc}
\synctex=-1
\usepackage{amsthm}

\makeatletter
%%%%%%%%%%%%%%%%%%%%%%%%%%%%%% Textclass specific LaTeX commands.
\newlength{\lyxlabelwidth}      % auxiliary length 
\theoremstyle{plain}
\newtheorem{thm}{\protect\theoremname}
\theoremstyle{plain}
\newtheorem{prop}[thm]{\protect\propositionname}
\theoremstyle{remark}
\newtheorem*{claim*}{\protect\claimname}

%%%%%%%%%%%%%%%%%%%%%%%%%%%%%% User specified LaTeX commands.
\date{}
\usepackage{commath}

\makeatother

\usepackage{babel}
\providecommand{\claimname}{Claim}
\providecommand{\propositionname}{Proposition}
\providecommand{\theoremname}{Theorem}

\begin{document}
I recommend we consider the term \emph{framed lattice} rather than
\emph{framed binary lattice} (i.e., ``binary'' is implied by ``framed'').
That might be good on its own but also is parallel to what I have
below.

I don't think the terminology I'm introducing is important enough
to define separately because it's only used in this one result. Therefore
I have written the terminology into the proof. I've also relegated
a proposition to a mere claim inside the proof. And what used to be
a theorem is here a proposition. We can discuss this, but my general
feeling is that a theorem is something that people want to know outside
the paper. That is, it has value outside the context of what we're
doing at that particular time in the paper.

\bigskip{}
\bigskip{}

\begin{prop}
For all positive integers $m$ and $n$, the $\left(0,0\right)$ entry
of $A_{n}^{m}$ equals $\sum v\left(\ell\right)$, where the sum is
over all $\left(m,n\right)$ framed binary lattices $\ell$ {[}{[}or,
using my suggestion above, ``\ldots{} over all $\left(m,n\right)$
framed lattices $\ell$''{]}{]}.
\end{prop}

\begin{proof}
Fix $n$. We first introduce some terminology. For all positive integers
$m$, we say a \emph{height~$m$ strip} is an $\left(m,n\right)$
binary lattice in which all vertices in the outside columns (i.e.,
column~$0$ and column~$n$) are labeled zero. {[}{[}Do we say ``vertices
are zero'' or ``vertices are labeled zero''? I'm not sure I care,
except that we should be consistent.{]}{]} Then for a height~$m$
strip $\ell$ and a nonnegative integer $i<2^{n-1}$, we say a particular
row of $\ell$ is \emph{determined by $i$} if that row is equal to
$0\ \beta_{n-1}\left(i\right)\ 0$. Note that an $\left(m,n\right)$
framed {[}{[}binary{]}{]} lattice is exactly a height~$m$ strip
in which both the top and bottom rows are determined by $0$. Thus,
to prove our proposition it suffices to show the following.
\begin{claim*}
For all positive integers $m$ and nonnegative integers $i,j<2^{n-1}$,
the $\left(i,j\right)$ entry of $A_{n}^{m}$ equals $\sum v\left(\ell\right)$,
where the sum is over all height~$m$ strips $\ell$ whose first
row is determined by $i$ and whose last row is determined by $j$.
\end{claim*}
To prove the claim, first observe that for each $i$ and $j$, there
exists a unique height~$1$ strip whose first row is determined by
$i$ and whose last row is determined by $j$. By Proposition~4.1,
$v$ applied to this strip is equal to the $\left(i,j\right)$ entry
of $A_{n}$. This proves our claim when $m=1$. We continue by induction,
so for some $m\ge2$, assume the claim for $m-1$.

Fix nonnegative integers $i,j<2^{n-1}$, and consider some nonnegative
integer $k<2^{n-1}$. As in the $m=1$ case, Proposition~4.1 tells
us that the $\left(i,k\right)$ entry of $A_{n}$ equals $v$ applied
to the unique height~$1$ strip with first and last rows determined
respectively by $i$ and by $k$. By our induction hypothesis, the
$\left(k,j\right)$ entry of $A_{n}^{m-1}$ equals $\sum v\left(\ell\right)$,
where the sum is over all height~$m-1$ strips $\ell$ whose first
and last rows are determined respectively by $k$ and by $j$. Consider
applying $v$ both to the height~$1$ strip indicated above and to
one of the height~$m-1$ strips $\ell$ in this sum. As in Figure~(picture
below), by the definition of $v$, the product of these two numbers
is equal to $v$ applied to the height~$m$ strip formed by overlapping
the bottom row of the height~$1$ strip to the top row of $\ell$.
\[
v\left(\begin{array}{ccc}
0 & \beta_{n-1}\left(i\right) & 0\\
0 & \beta_{n-1}\left(k\right) & 0
\end{array}\right)v\left(\begin{array}{ccc}
0 & \beta_{n-1}\left(k\right) & 0\\
\vdots & \vdots & \vdots\\
0 & \beta_{n-1}\left(j\right) & 0
\end{array}\right)=v\left(\begin{array}{ccc}
0 & \beta_{n-1}\left(i\right) & 0\\
0 & \beta_{n-1}\left(k\right) & 0\\
\vdots & \vdots & \vdots\\
0 & \beta_{n-1}\left(j\right) & 0
\end{array}\right)
\]

From this we see that the product of the $\left(i,k\right)$ entry
of $A_{n}$ and the $\left(k,j\right)$ entry of $A_{n}^{m-1}$ is
equal to $\sum v\left(\ell\right)$, where the sum is now over all
height~$m$ strips $\ell$ whose first, second, and last rows are
determined respectively by $i$, by $k$, and by $j$. This tells
us, in turn, that when we additionally sum over all nonnegative $k<2^{n-1}$,
we can conclude that the $\left(i,j\right)$ entry of $A_{n}\cdot A_{n}^{m-1}=A_{n}^{m}$
is, as desired, the sum $\sum v\left(\ell\right)$ where the sum is
over all height~$m$ strips $\ell$ whose first and last rows are
determined respectively by $i$ and by $j$.
\end{proof}

\end{document}
